% ============ 第4章 §4.5 非退化振荡积分算子(原书 pp.336-338) ============
\section{非退化振荡积分算子}\label{sec:4.5}
\index{振荡积分算子}
\index{非退化条件}


作为上一节的自然推广,本节讨论非退化振荡积分算子.

设 $a(x, y) \in C_{0}^{\infty}(\mathbb{R}^{n}\times\mathbb{R}^{n})$.对实值函数 $\phi(x, y) \in C^{\infty}(\mathbb{R}^{n}\times\mathbb{R}^{n})$,考虑振荡积分算子
\begin{equation*}
T_{\lambda}f(x) := \int_{\mathbb{R}^{n}}e^{i\lambda\phi(x, y)}a(x, y)f(y)\dd y \quad (\lambda > 1)
\end{equation*}
的估计.当然,目标是数量地把握 $\lambda \to \infty$ 时的衰减.以下对 $\phi(x, y)$ 假设非退化条件
\begin{equation}
\det\left((\partial^{2}/\partial x_{j}\partial y_{k})\phi(x, y)\right) \neq 0.\tag{4.5.1}
\label{eq:4.5.1}\end{equation}

\begin{theorem}\label{thm:4.5.1}
若 $\phi(x, y)$ 在 $a(x, y)$ 的支集上满足 \eqref{eq:4.5.1},则
\begin{equation}
\left\|\int_{\mathbb{R}^{n}}e^{i\lambda\phi(x, y)}a(x, y)f(y)\dd y\right\|_{L^{2}(\mathbb{R}^{n})} \leq C\lambda^{-n/2}\|f\|_{L^{2}(\mathbb{R}^{n})}\tag{4.5.2}
\label{eq:4.5.2}\end{equation}
成立.
\end{theorem}

显然,
\begin{equation}
\left\|\int_{\mathbb{R}^{n}}e^{i\lambda\phi(x, y)}a(x, y)f(y)\dd y\right\|_{L^{\infty}(\mathbb{R}^{n})} \leq C\|f\|_{L^{1}(\mathbb{R}^{n})}.\tag{4.5.3}
\label{eq:4.5.3}\end{equation}
由 M. Riesz 插值定理插值 \eqref{eq:4.5.2},\eqref{eq:4.5.3},得到下面的系.

\begin{corollary}\label{cor:4.5.2}
设 $1 \leq p \leq 2$,$p' = p/(p-1)$,则
\begin{equation}
\|T_{\lambda}f\|_{L^{p'}(\mathbb{R}^{n})} \leq C\lambda^{-n/p'}\|f\|_{L^{p}(\mathbb{R}^{n})}\tag{4.5.4}
\label{eq:4.5.4}\end{equation}
成立.
\end{corollary}

若 $\phi(x, y) = -2\pi ix\cdot y$,则这满足 \eqref{eq:4.5.1} 的条件.变换变量,由 \eqref{eq:4.5.2} 可得
\begin{equation*}
\left\|\int_{\mathbb{R}^{n}}e^{-2\pi i x\cdot y}a(x/\sqrt{\lambda}, y/\sqrt{\lambda})f(y)\dd y\right\|_{L^{2}(\mathbb{R}^{n})} \leq C\|f\|_{L^{2}(\mathbb{R}^{n})}.
\end{equation*}
选取 $a(x, y)$ 在 $(x, y) = (0, 0)$ 处取 1,令 $\lambda \to \infty$,即得 Fourier 变换的 $L^{2}$ 有界性.

\begin{proof}[定理~\ref{thm:4.5.1} 的证明]
由 \eqref{eq:4.5.1} 的条件与 Taylor 定理,可以选取 $c > 0$,使得当 $|y - z|$ 充分小时
\begin{equation}
|\nabla_{x}[\phi(x, y) - \phi(x, z)]| \geq c|y - z|\tag{4.5.5}
\label{eq:4.5.5}\end{equation}
成立.必要时将 $a(x, y)$ 的支集用光滑的有限个单位分解分解,可假设在 $a(x, y)$ 的支集上 \eqref{eq:4.5.5} 成立.

首先改写为
\begin{equation}
\begin{aligned}
\|T_{\lambda}f\|_{2}^{2} & = \int\int K_{\lambda}(y, z)f(y)\overline{f(z)}\dd y\dd z, \\
K_{\lambda}(y, z) & := \int_{\mathbb{R}^{n}}e^{i\lambda[\phi(x, y) - \phi(x, z)]}a(x, y)\overline{a(x, z)}\dd x.
\end{aligned}
\end{equation}
考虑 \eqref{eq:4.5.5},与定理~\ref{thm:4.4.1} 的证明同样,可估计为
\begin{equation*}
|K_{\lambda}(y, z)| \leq C_{N}(1 + \lambda|y - z|)^{-N} \quad (N = 1, 2, \ldots).
\end{equation*}
取 $N = n + 1$,
\begin{equation*}
\int_{\mathbb{R}^{n}}(1 + \lambda|y|)^{-(n+1)}\dd y = C\lambda^{-n},
\end{equation*}
由 Young 不等式,
\begin{equation*}
\left\|\int(1 + \lambda|y - z|)^{-(n+1)}|f(y)|\dd y\right\|_{2} \leq C\lambda^{-n}\|f\|_{2}.
\end{equation*}
故由 Schwarz 不等式,
\begin{equation*}
\|T_{\lambda}f\|_{2}^{2} \leq \int\int(1 + \lambda|y - z|)^{-(n+1)}|f(y)||f(z)|\dd y\dd z \leq C\lambda^{-n}\|f\|_{2}^{2}.
\end{equation*}
到达目标.
\end{proof}
